\documentclass[11pt]{article}
\usepackage[margin=1in]{geometry}
\usepackage{amsmath,amssymb,amsthm,mathtools}
\usepackage[hidelinks]{hyperref}
\usepackage{microtype}
\usepackage[T1]{fontenc}

\newtheorem{theorem}{Theorem}
\newtheorem{lemma}{Lemma}
\newcommand{\Av}{\operatorname{Av}}
\newcommand{\SEQ}{\operatorname{SEQ}}

\title{The Stanley--Wilf Limit\\
  \large A short proof and a Lean 4 formalization}
\author{Xiang Huang}
\date{September 2026}

\begin{document}
\maketitle

\begin{abstract}
For every fixed permutation pattern $\tau$, we prove that
$|\Av_n(\tau)|^{1/n}$ has a finite limit.  The structural step is the
symbolic specification of the avoidance class as a sequence of
indecomposable permutations, which gives supermultiplicativity.  Fekete's
lemma then gives existence of the limit, while the Marcus--Tardos theorem
gives finiteness.  We also describe a kernel-checked Lean 4 formalization in
which the Marcus--Tardos extremal estimate and the intervening matrix-counting
argument are themselves formalized.
\end{abstract}

Let $\sigma\oplus\pi$ denote the direct sum of two permutations and
$\sigma\ominus\pi$ their skew sum.  A nonempty permutation is
\emph{sum-indecomposable} (respectively \emph{skew-indecomposable}) if it has
no nontrivial decomposition using $\oplus$ (respectively $\ominus$).

\begin{lemma}\label{lem:dichotomy}
A permutation of length at least two is sum-indecomposable or
skew-indecomposable.
\end{lemma}

\begin{proof}
If $\tau=\alpha\oplus\beta$ is a nontrivial direct sum, then the entry $1$
occurs before the largest entry.  If
$\tau=\gamma\ominus\delta$ is a nontrivial skew sum, the largest entry occurs
before $1$.  Both decompositions therefore cannot exist.
\end{proof}

\begin{lemma}\label{lem:closure}
If $\tau$ is sum-indecomposable, then $\Av(\tau)$ is closed under direct sum.
If $\tau$ is skew-indecomposable, it is closed under skew sum.
\end{lemma}

\begin{proof}
We prove the direct-sum case.  Suppose that $\sigma$ and $\pi$ avoid $\tau$.
An occurrence of $\tau$ in $\sigma\oplus\pi$ either lies in one block, which
contradicts avoidance by that block, or meets both blocks.  In the latter
case every selected value from the first block is below every selected value
from the second block.  The occurrence therefore expresses $\tau$ as a
nontrivial direct sum, a contradiction.  The skew case is identical with the
inequalities reversed.
\end{proof}

\begin{theorem}[Stanley--Wilf]
For every permutation pattern $\tau$, the limit
\[
  \lim_{n\to\infty}|\Av_n(\tau)|^{1/n}
\]
exists and is finite.
\end{theorem}

\begin{proof}
Write $c_n=|\Av_n(\tau)|$.  Patterns of length at most one are immediate: all
positive-size avoidance classes are empty, so the limit is zero.  Assume
$|\tau|\geq2$.

By Lemma~\ref{lem:dichotomy}, choose direct sum or skew sum so that $\tau$ is
indecomposable for that operation.  Lemma~\ref{lem:closure} shows that
$\Av(\tau)$ is closed under the chosen operation.  Let $\mathcal I$ be its
nonempty indecomposable members.  The unique decomposition into
indecomposable components gives the symbolic specification
\[
  \Av(\tau)=\SEQ(\mathcal I),
  \qquad C(z)=\frac{1}{1-I(z)}.
\]
In particular, concatenating the component sequences of objects of sizes
$m$ and $n$ gives an injection
\[
  \Av_m(\tau)\times\Av_n(\tau)\hookrightarrow\Av_{m+n}(\tau).
\]
Thus $c_{m+n}\geq c_mc_n$.

The singleton permutation avoids $\tau$, so repeated sums give $c_n\geq1$.
Hence $\log c_n$ is superadditive.  Fekete's lemma yields
\[
  \lim_{n\to\infty}\frac{\log c_n}{n}
    =\sup_{n\geq1}\frac{\log c_n}{n},
\]
initially allowing $+\infty$.  Exponentiating gives
\[
  \lim_{n\to\infty}c_n^{1/n}=\sup_{n\geq1}c_n^{1/n}.
\]
Finally, the Marcus--Tardos theorem supplies a constant $K_\tau<\infty$ such
that $c_n\leq K_\tau^n$ for all $n$.  The displayed supremum is therefore at
most $K_\tau$.
\end{proof}

\section*{Formalization}

The Lean development formalizes both inputs rather than postulating the
exponential bound.  For a pattern of length $k$, the formalized
Marcus--Tardos extremal estimate uses the explicit linear constant
\[
  M_k=2k^4\binom{k^2}{k}.
\]
The permutation-matrix embedding, a $2\times2$ block contraction, and the
fifteen possible nonempty block masks give Klazar's recurrence.  A dyadic
argument then proves the explicit bound
\[
  c_n\leq\bigl(2\cdot15^{2M_k}\bigr)^n.
\]
The public Lean theorem has no mathematical hypotheses:
\begin{verbatim}
theorem stanleyWilf (tau : Perm k) : GrowthTarget tau
\end{verbatim}
where \texttt{GrowthTarget tau} states that the real sequence
$c_n^{1/n}$ tends to a finite nonnegative real number.

The full project builds with Lean~4.29.0.  A transitive
\texttt{\#print axioms} audit of 48 declarations reports only
\texttt{propext}, \texttt{Classical.choice}, and \texttt{Quot.sound}.
The source and verification scripts are available at
\href{https://github.com/xiangyazi24/stanley-wilf}
{github.com/xiangyazi24/stanley-wilf}.

The extremal-matrix layer is an attributed Lean~4.29 adaptation of Ya\"el
Dillies and contributors' Apache-2.0 \texttt{ForbiddenMatrix} development.
The permutation bridge, counting recurrence, dyadic bound, symbolic
decomposition, and final theorem are assembled in this repository.

\begin{thebibliography}{9}
\bibitem{Arratia}
R. Arratia,
\emph{On the Stanley--Wilf conjecture for the number of permutations avoiding
a given pattern}, Electron. J. Combin. 6 (1999), N1.

\bibitem{MarcusTardos}
A. Marcus and G. Tardos,
\emph{Excluded permutation matrices and the Stanley--Wilf conjecture},
J. Combin. Theory Ser. A 107 (2004), 153--160.

\bibitem{FlajoletSedgewick}
P. Flajolet and R. Sedgewick,
\emph{Analytic Combinatorics}, Cambridge University Press, 2009.
\end{thebibliography}

\end{document}
